\documentclass[11pt]{article}

\usepackage[margin=1in]{geometry}
\usepackage{amsmath,amssymb}
\usepackage{xcolor}
\usepackage{graphicx}
\usepackage{hyperref}
\usepackage{booktabs}
\usepackage{listings}

\definecolor{codegreen}{rgb}{0,0.6,0}
\definecolor{codegray}{rgb}{0.5,0.5,0.5}
\definecolor{codepurple}{rgb}{0.58,0,0.82}
\definecolor{backcolour}{rgb}{0.95,0.95,0.92}

\lstdefinestyle{mystyle}{
  backgroundcolor=\color{backcolour},
  commentstyle=\color{codegreen},
  keywordstyle=\color{magenta},
  numberstyle=\tiny\color{codegray},
  stringstyle=\color{codepurple},
  basicstyle=\ttfamily\footnotesize,
  breakatwhitespace=false,
  breaklines=true,
  captionpos=b,
  keepspaces=true,
  numbers=left,
  numbersep=5pt,
  showspaces=false,
  showstringspaces=false,
  showtabs=false,
  tabsize=2
}
\lstset{style=mystyle}

\title{Training Run v2: Multi-Trace Augmentation, Throughput Debugging, and Reward~A}
\author{Pulak Mehrotra}
\date{March 28--30, 2026}

\begin{document}

\maketitle

\begin{abstract}
This report documents the first end-to-end training runs of the Octopus
CXL memory-pooling RL agent using the full 10-trace augmented dataset
(\texttt{--multi-trace}, \texttt{--augmentation}, $N_\text{envs}=32$).
Three system-level bugs were identified and fixed: a copy-on-write OOM
during \texttt{SubprocVecEnv} fork, a pickle import-lock deadlock in
parallel trace loading, and a \texttt{KeyError: 192} crash in
\texttt{PoolingSavingsCallback} caused by building \texttt{mhd\_to\_hosts}
from the base pod matrix rather than the full block-diagonal
\texttt{expanded\_M}.
Throughput profiling identified eval callbacks (not gradient updates or
environment steps) as the true bottleneck, and a \texttt{train\_freq=8},
\texttt{eval\_freq=100000} configuration restores $\approx$2,400~fps.
The original reward (\texttt{current}) showed near-zero policy improvement;
Reward~A with a 50-dim observation space is the active variant under test.
\end{abstract}

% ─────────────────────────────────────────────────────────────────────────────
\section{Run Configuration}
% ─────────────────────────────────────────────────────────────────────────────

\begin{lstlisting}[language=bash]
python scripts/train_rl.py \
  --run-id aug_all_traces \
  --augmentation \
  --multi-trace \
  --aug-traces AMS20PrdApp19-tround BLAPrdApp19-troundgrt5m \
               BN9PrdApp18-troundgrt5m DSM08PrdApp05-troundgrt5m \
               DUB24PrdApp09-troundgrt5m LON23PrdApp01-troundgrt5m \
               LVL01PrdApp05-troundgrt5m SG2PrdApp35-troundgrt5m \
               SYD21PrdApp07-troundgrt5m YTO21PrdApp05-troundgrt5m \
  --wandb --n-envs 32 --total-timesteps 500000
\end{lstlisting}

\noindent Key parameters: reward variant = \texttt{current}, $\lambda_\text{var}=0.5$,
batch size~=~256, $\gamma=0.999$, \texttt{train\_freq}=1,
\texttt{gradient\_steps}=1, network~=~[256,~256] MLP,
eval trace = \texttt{LON23PrdApp01-troundgrt5m}.

% ─────────────────────────────────────────────────────────────────────────────
\section{System Fixes Required}
% ─────────────────────────────────────────────────────────────────────────────

\subsection{Copy-on-Write OOM During Worker Fork}

The original run crashed silently after loading all 10 traces.
The OOM killer terminated the process with no Python traceback.

\textbf{Root cause.} Python's reference counting increments
\texttt{ob\_refcnt} on every read of a Python object, which marks the
memory page as dirty and triggers a copy-on-write copy in each forked
subprocess. With 10 traces loaded as Python object graphs
($\sim$53~GB in memory) and 32 workers each copying $\sim$7~GB of
CoW pages, peak memory reached $\sim$277~GB on a 192~GB machine.

\textbf{Fix.} Trace data is now converted to \texttt{TraceArrays}
(7 numpy arrays: \texttt{vm\_start}, \texttt{vm\_end}, \texttt{vm\_mem},
\texttt{node\_ids}, \texttt{node\_dram}, \texttt{node\_offsets},
\texttt{vm\_ptrs}) via \texttt{octopus.data.to\_arrays()} before the
\texttt{SubprocVecEnv} fork. Reading numpy C-heap buffers does not
touch Python refcounts, so no CoW pages are generated.
Total pool footprint dropped from $\sim$53~GB to $\sim$300~MB.

\subsection{Pickle Import-Lock Deadlock in Parallel Loading}

A \texttt{ThreadPoolExecutor} was used to load all 10 traces
simultaneously. This deadlocked because \texttt{pickle.Unpickler.find\_class}
calls \texttt{\_\_import\_\_}, which acquires Python's per-module import
lock. Ten threads all calling \texttt{\_\_import\_\_} simultaneously
caused a classic import-lock deadlock.

\textbf{Fix.} Reverted to sequential loading. Startup time is
$\approx$10.5~min, dominated by I/O-bound pickle deserialization and
the \texttt{to\_arrays()} conversion.

% ─────────────────────────────────────────────────────────────────────────────
\section{Startup Time Breakdown}
% ─────────────────────────────────────────────────────────────────────────────

\begin{table}[h]
\centering
\begin{tabular}{lrr}
\toprule
Phase & Wall time & Cumulative \\
\midrule
W\&B init                      & 0:00 -- 0:09:30 & 9 min 30 s \\
Training trace load + topology  & 0:09:30 -- 0:10:00 & 30 s \\
10-trace pool load (sequential) & 0:10:00 -- 0:18:31 & 8 min 31 s \\
SubprocVecEnv fork (32 workers) & 0:18:31 -- 0:20:34 & 2 min 03 s \\
Eval trace load                 & 0:20:34 -- 0:20:56 & 22 s \\
\midrule
\textbf{Total to first step}    & & \textbf{$\approx$20 min 56 s} \\
\bottomrule
\end{tabular}
\caption{Startup phases. Excluding W\&B, data loading accounts for
$\approx$11 min 26 s.}
\end{table}

The 10-trace pool load scales linearly with trace size. The largest
traces (SG2: 1.7~GB, BLA: 1.6~GB, BN9: 1.5~GB on disk) dominate,
taking 60--80~s each. \texttt{to\_arrays()} conversion adds a fixed
per-trace overhead of $\approx$5--15~s (tight Python loop over VM
objects to extract int64 timestamps and float32 memory values).

% ─────────────────────────────────────────────────────────────────────────────
\section{Training Throughput}
% ─────────────────────────────────────────────────────────────────────────────

\subsection{Observed vs.\ Expected}

\begin{table}[h]
\centering
\begin{tabular}{lrrr}
\toprule
Checkpoint & \texttt{time\_elapsed} (s) & \texttt{total\_timesteps} & fps \\
\midrule
115,296  & 1,431 & 115,296 & 80 \\
206,528  & 2,837 & 206,528 & 72 \\
215,488  & 2,850 & 215,488 & 75 \\
\midrule
\textbf{Average} & & & \textbf{75} \\
\bottomrule
\end{tabular}
\caption{Observed training throughput. fps = \texttt{total\_timesteps} /
\texttt{time\_elapsed}.}
\end{table}

The benchmark from training-v1 with $N=32$ SubprocVecEnv was
\textbf{1,221~fps}. The current run achieves \textbf{75~fps} --- a
\textbf{16$\times$ regression}.

\subsection{Projected Wall Times}

\begin{table}[h]
\centering
\begin{tabular}{lrr}
\toprule
Total timesteps & Training time & + Startup ($\approx$11 min) \\
\midrule
500k   & 1.85 hr & $\approx$2.0 hr \\
1M     & 3.7 hr  & $\approx$3.9 hr \\
2M     & 7.4 hr  & $\approx$7.6 hr \\
5M     & 18.5 hr & $\approx$18.7 hr \\
\bottomrule
\end{tabular}
\caption{Projected run times at 75~fps. Overnight budget $\approx$8~hr
$\Rightarrow$ 2M timesteps feasible.}
\end{table}

\subsection{Root Cause: Eval Callbacks, Not Gradient Updates}
\label{sec:rootcause}

Two profiling experiments identified the true bottleneck.

\textbf{Experiment 1 — environment-only profiling} (no gradient updates,
\texttt{scripts/profile\_throughput.py}):

\begin{table}[h]
\centering
\begin{tabular}{lrr}
\toprule
Configuration & fps & vs.\ baseline \\
\midrule
Single trace, no aug, no multi-trace & 2,806 & 1.0$\times$ \\
+ Multi-trace only                    & 2,611 & 0.93$\times$ \\
+ Augmentation only                   & 2,760 & 0.98$\times$ \\
+ Both (current training config)      & 2,565 & 0.91$\times$ \\
\bottomrule
\end{tabular}
\caption{Environment throughput without gradient updates.
Augmentation and multi-trace have negligible impact ($<$10\%).
All four run $\gg$75~fps — the env is not the bottleneck.}
\end{table}

\textbf{Experiment 2 — \texttt{train\_freq} sweep with gradient updates}
(\texttt{scripts/smoke\_train\_freq.py}), $N=32$ envs, 3k steps, no callbacks:

\begin{table}[h]
\centering
\begin{tabular}{rrrrrr}
\toprule
\texttt{train\_freq} & fps & vs.\ tf=1 & grad updates & grad mean & collect mean \\
\midrule
 1 & 1,031 & 1.0$\times$ & 79 & 24 ms & ~5 ms \\
 4 & 1,790 & 1.7$\times$ & 21 & 24 ms & 23 ms \\
 8 & 2,454 & 2.4$\times$ & 11 & 19 ms & 33 ms \\
16 & 2,715 & 2.6$\times$ &  6 & 17 ms & 65 ms \\
32 & 2,484 & 2.4$\times$ &  3 & 24 ms & 167 ms \\
\bottomrule
\end{tabular}
\caption{Each gradient update costs $\approx$20--25~ms. The system
without callbacks reaches 1,000--2,700~fps depending on \texttt{train\_freq}.
\texttt{train\_freq}=8 is the sweet spot: 2.4$\times$ speedup with 2$\times$
more updates than tf=16 for better sample efficiency.}
\end{table}

\textbf{Conclusion.} The env+gradient system achieves
$\approx$1,000~fps at \texttt{train\_freq}=1 without callbacks.
Real training achieves 75~fps.
The gap — $\approx$266~seconds per 20k-step eval window — is entirely
attributable to eval callbacks:

\begin{itemize}
  \item \textbf{\texttt{PoolingSavingsCallback}} (dominant cost): runs
        \texttt{pooling\_simulation} for 10 full inference episodes every
        \texttt{--eval-freq=20000} steps. Each inference episode is
        $\sim$7,000 serial steps. Total: $\sim$70,000 inference steps
        per eval window, amortised as $\approx$220~ms per training step.
  \item \textbf{\texttt{EvalCallback}}: 3 eval episodes $\times$~7,000 steps
        $\approx$21,000 inference steps. Amortised $\approx$65~ms per
        training step.
\end{itemize}

\subsection{Are the Callbacks Necessary?}

\begin{table}[h]
\centering
\begin{tabular}{llll}
\toprule
Callback & Affects training? & Cost & Verdict \\
\midrule
\texttt{EvalCallback}          & \textbf{Yes} — saves \texttt{best\_model.zip} & Medium & Keep, raise freq \\
\texttt{PoolingSavingsCallback} & No — diagnostic only                          & High   & Raise freq or drop \\
\texttt{CheckpointCallback}    & No — periodic snapshots                       & Low    & Keep \\
\texttt{SACDiagnosticsCallback} & No — logs Q/entropy                          & Negligible & Keep \\
\texttt{WandbCallback}         & No — artifact upload                          & Low    & Keep \\
\bottomrule
\end{tabular}
\caption{\texttt{PoolingSavingsCallback} does \textbf{not} save checkpoints or
influence the gradient signal. It is purely diagnostic and accounts for
$\sim$80\% of callback overhead.}
\end{table}

\texttt{EvalCallback} is worth retaining because it is the only mechanism
that saves the best policy during training --- critical if the reward
improves then regresses. However, both callbacks should run far less frequently.

\subsection{Projected Wall Times After Fix}

Recommended config: \texttt{--train-freq 8 --eval-freq 100000}.

\begin{table}[h]
\centering
\begin{tabular}{lrr}
\toprule
Total timesteps & Estimated time & Notes \\
\midrule
500k  & $\sim$12 min  & \\
1M    & $\sim$24 min  & \\
2M    & $\sim$48 min  & Overnight budget \\
5M    & $\sim$2 hr    & \\
\bottomrule
\end{tabular}
\caption{Projected at $\approx$2,400~fps (tf=8, eval-freq=100k).
Startup adds $\approx$11~min.}
\end{table}

% ─────────────────────────────────────────────────────────────────────────────
\section{Learning Signal Analysis}
% ─────────────────────────────────────────────────────────────────────────────

\subsection{Pooling Savings}

\begin{table}[h]
\centering
\begin{tabular}{rrr}
\toprule
Timesteps & \texttt{pooling\_savings\_mean} & \texttt{pooling\_savings\_std} \\
\midrule
20k  & 0.0048 & 0.0275 \\
40k  & 0.0128 & 0.0164 \\
60k  & 0.0156 & 0.0131 \\
80k  & 0.0140 & 0.0128 \\
100k & 0.0139 & 0.0119 \\
120k & 0.0131 & 0.0148 \\
140k & 0.0139 & 0.0151 \\
160k & 0.0101 & 0.0132 \\
180k & 0.0068 & 0.0131 \\
200k & 0.0139 & 0.0090 \\
\bottomrule
\end{tabular}
\caption{\texttt{pooling\_savings\_mean} = fraction by which RL outperforms
the greedy baseline on the eval trace (LON23). Standard deviation across
10 eval episodes.}
\end{table}

The policy achieves \textbf{$\approx$1.4\% savings over greedy} at
peak, but this improvement is \textbf{not monotonically increasing} and
shows no upward trend after 60k steps. The savings are within one
standard deviation of zero at most checkpoints, suggesting the policy
has not found a meaningfully better strategy than the greedy baseline.

\subsection{SAC Internals}

Several internal SAC metrics signal problematic learning dynamics:

\begin{itemize}
  \item \textbf{Entropy collapse}: \texttt{ent\_coef} decays
        monotonically from 0.837 (20k steps) to 0.155 (160k steps).
        The policy is becoming deterministic, but without finding a
        better optimum.
  \item \textbf{Unchecked Q-value growth}: \texttt{q1\_mean} rises
        from 12.5 to $>$70 over the run. This is consistent with
        the Q-function tracking cumulative negative reward without a
        bound, but at a rate suggesting value overestimation.
  \item \textbf{Negative training-env savings}: \texttt{env/pooling\_savings}
        is negative throughout ($-0.10$ to $-0.44$), meaning the
        policy performs \textit{worse} than greedy on the training
        traces --- it only marginally outperforms greedy on the
        held-out eval trace.
  \item \textbf{Critic loss converging}: \texttt{critic\_loss} drops
        from 0.074 to 0.004, indicating the value function has
        converged, but to a poor policy.
\end{itemize}

\subsection{Why Reward ``current'' Fails}

The reward $R = -(\Delta\text{peak})/\text{fair\_share} - \lambda\cdot\text{Var}(\mathbf{c}/\text{fair\_share})$
has three structural problems:

\begin{enumerate}
  \item \textbf{Sparse signal}: only the single worst MPD contributes
        to the peak-delta term. Most decisions receive near-zero gradient.
  \item \textbf{Myopic}: no lookahead. The agent cannot anticipate
        departures or future arrivals, so it cannot balance load
        over time.
  \item \textbf{Scale sensitivity}: the fair-share normalization depends
        on pod-level DRAM, which varies per episode under multi-trace
        augmentation (5--7$\times$ scale range). The reward magnitude
        is inconsistent across episodes, harming value function stability.
\end{enumerate}

% ─────────────────────────────────────────────────────────────────────────────
\section{First Reward~A Run: Crash and Fix}
\label{sec:rewardA-crash}
% ─────────────────────────────────────────────────────────────────────────────

Armed with the throughput fixes (\texttt{train\_freq=8}, \texttt{eval\_freq=100000}),
the first Reward~A run (\texttt{--run-id aug\_v2\_rewardA}) was launched.

\subsection{Run Configuration}

\begin{lstlisting}[language=bash]
python scripts/train_rl.py \
  --run-id aug_v2_rewardA \
  --reward-variant A \
  --augmentation --multi-trace \
  --aug-traces AMS20PrdApp19-tround BLAPrdApp19-troundgrt5m \
               BN9PrdApp18-troundgrt5m DSM08PrdApp05-troundgrt5m \
               DUB24PrdApp09-troundgrt5m LON23PrdApp01-troundgrt5m \
               LVL01PrdApp05-troundgrt5m SG2PrdApp35-troundgrt5m \
               SYD21PrdApp07-troundgrt5m YTO21PrdApp05-troundgrt5m \
  --n-envs 32 --total-timesteps 2000000 \
  --train-freq 8 --eval-freq 100000
\end{lstlisting}

\subsection{Early Signal at 100k Steps}

The first eval completed successfully at 100k steps:

\begin{table}[h]
\centering
\begin{tabular}{lr}
\toprule
Metric & Value \\
\midrule
\texttt{pooling\_ratio}      & 1.05 \\
\texttt{pooling\_savings}    & $-$0.053 (5.3\% worse than greedy) \\
\texttt{mean\_reward}        & $-$25.1 \\
\texttt{episode\_length}     & 7,334 \\
\texttt{critic\_loss}        & 0.112 \\
\texttt{ent\_coef}           & 0.891 \\
\texttt{n\_updates}          & 387 \\
\bottomrule
\end{tabular}
\caption{Eval metrics at 100k steps. High entropy (0.891) and elevated critic loss
indicate the policy is still in early exploration. Negative savings at 100k is
expected for Reward~A --- the 50-dim observation requires more experience before
the value function stabilises.}
\end{table}

\subsection{Crash: \texttt{KeyError: 192} in \texttt{PoolingSavingsCallback}}

The run crashed immediately after the 100k-step eval (first checkpoint saved),
during the next \texttt{PoolingSavingsCallback} invocation:

\begin{lstlisting}
File "scripts/evaluate.py", line 344, in _rl_alloc_cb
    P_j = float(sum(host_cxl_load[h] for h in mhd_to_hosts[mhd]))
KeyError: 192
\end{lstlisting}

\textbf{Root cause.} \texttt{expand\_M\_to\_all\_nodes} assembles a
block-diagonal topology matrix where pod~0 occupies MPD columns 0--191,
pod~1 occupies columns 192--383, and so on.
The \texttt{mhd\_to\_hosts} dict was precomputed once from the base
per-pod matrix~$M$ (192 entries, keys 0--191) and stored in the callback
at startup.
When the callback evaluated a pod sampled from block~1 or higher,
\texttt{mhd\_list} contained global column indices $\geq$192,
which were absent from \texttt{mhd\_to\_hosts}.

The first eval (100k steps) happened to sample only pod~0 in all 10
iterations (seed~99999--100008 all landed there), masking the bug.
The second eval (200k steps) sampled at least one pod with index~$\geq$1,
triggering the crash.

\textbf{Fix.} \texttt{PoolingSavingsCallback.\_run\_eval()} now rebuilds
\texttt{mhd\_to\_hosts} and \texttt{Q\_j} from \texttt{expanded\_M}
inside each iteration, immediately after \texttt{expand\_M\_to\_all\_nodes}
returns.  The rebuilt structures cover all $\text{num\_pods} \times 192$
global MPD indices:

\begin{lstlisting}[language=Python]
node_to_M, expanded_M = expand_M_to_all_nodes(self._M, pod_to_nodes)
if self._obs_variant != "current":
    num_mhd_exp = len(expanded_M[0])
    mhd_to_hosts_exp = {j: [] for j in range(num_mhd_exp)}
    h_to_mhds = {}
    for h, row in enumerate(expanded_M):
        h_to_mhds[h] = [j for j, v in enumerate(row) if v]
        for j in h_to_mhds[h]:
            mhd_to_hosts_exp[j].append(h)
    Q_j_exp = np.zeros(num_mhd_exp, dtype=np.float32)
    for j in range(num_mhd_exp):
        hosts = mhd_to_hosts_exp[j]
        if hosts:
            inv_deg = [1.0/len(h_to_mhds[h]) for h in hosts if h_to_mhds[h]]
            if inv_deg:
                Q_j_exp[j] = float(np.mean(inv_deg))
\end{lstlisting}

The pre-built \texttt{self.\_mhd\_to\_hosts} (base-pod scope) is no
longer passed to \texttt{make\_rl\_alloc\_cb} for obs variants other
than \texttt{current}.

% ─────────────────────────────────────────────────────────────────────────────
\section{Run~v3: Callback Speed Regression and Inference Analysis}
\label{sec:run3}
% ─────────────────────────────────────────────────────────────────────────────

\subsection{Observed Throughput}

With the \texttt{KeyError} fix in place, run~v3 (\texttt{aug\_v2\_rewardA})
trained at $\approx$43~fps overall --- worse than the 75~fps seen in the
original run with \texttt{obs\_variant=current}.
Between callback invocations the env+gradient rate was $\approx$860~fps,
confirming the fix was correct and the training loop itself was healthy.
Each \texttt{PoolingSavingsCallback} invocation was blocking for
$\approx$38--39~minutes:

\begin{center}
\begin{tabular}{lrr}
\toprule
Event & Wall time & \texttt{time\_elapsed} (s) \\
\midrule
Training start                       & 05:04:59 & 0 \\
\texttt{PoolingSavings @ 100k steps} & 05:44:49 & 2,395 \\
Training 104k $\to$ 182k (78k steps) & fast      & +88~s (860~fps) \\
\texttt{PoolingSavings @ 200k steps} & 06:25:23 & 4,839 \\
\bottomrule
\end{tabular}
\end{center}

100k environment steps at 2,400~fps takes $\approx$42~s.
The remaining $2{,}395 - 42 \approx 2{,}353$~s (39~min) per callback is
pure Python and inference overhead.

\subsection{Root Cause: D$_j$/S$_j$ Loop over Full Expanded Topology}

Switching from \texttt{obs\_variant=current} to \texttt{A} enables
\texttt{track\_vm\_allocs=True}, which activates the $D_j$/$S_j$ departure
lookahead computation in \texttt{\_rl\_alloc\_cb}:

\begin{lstlisting}[language=Python]
D_j = np.zeros(num_mhd, dtype=np.float64)   # num_mhd = num_pods * 192
for j in range(num_mhd):                     # ~11,904 iterations
    for dt, mem in mpd_vm_allocs[j]:
        ...
\end{lstlisting}

This loop ran for \textbf{every VM arrival event} across all pods.
With the LON23 eval trace spanning $\sim$10 pods
($\sim$100k total VM events across 10 simulation iterations),
the total Python iteration count was:

\[
100{,}000 \text{ events} \times 11{,}904 \text{ MPDs} \times 10 \text{ iters}
= 1.2\,\text{B iterations} \approx 38\text{ min at 10M\,iter/s}
\]

Only 8 of those 11,904 values were ever read (the \texttt{mhd\_list} for
the current host). The remaining 11,896 were dead computation.

\textbf{Fix.} Restrict $D_j$/$S_j$ to \texttt{mhd\_list} (max 8 MPDs):

\begin{lstlisting}[language=Python]
for k, mhd in enumerate(mhd_list):      # 8 iterations, not 11,904
    dj = sj = 0.0
    for dt, mem in mpd_vm_allocs[mhd]:
        if dt <= t_plus_W:
            dj += mem * (1.0 - (dt - tick) / W)
        else:
            sj += mem
\end{lstlisting}

Speedup on this loop: $11{,}904 / 8 \approx \mathbf{1{,}488\times}$.

\subsection{Inference Path Analysis}

With the $D_j$ fix, the remaining callback cost is dominated by
$\sim$400k--1M calls to \texttt{model.predict}.
Benchmarking the 50-dim $\to$ [256,\,256] $\to$ 8-dim MLP:

\begin{table}[h]
\centering
\begin{tabular}{lrr}
\toprule
Mode & ms/call & calls/s \\
\midrule
GPU forward + \texttt{.cpu().numpy()} & 0.167 & 5,999 \\
CPU forward (reused tensor)           & 1.156 & 864 \\
CPU forward (new tensor/call)         & 1.105 & 904 \\
\bottomrule
\end{tabular}
\caption{Single-sample inference cost for the 50-dim MLP.
GPU wins by $7\times$ even for batch size~1 because CUDA kernel launch
overhead (0.1\,ms) is lower than serial CPU GEMM for this model size.
Moving to CPU would make inference \textbf{slower}, not faster.
Batching is not possible: each allocation changes \texttt{cur\_cxl\_mem\_vec},
making the simulation strictly sequential.}
\end{table}

\noindent\textbf{Why async RL does not help here.}
Asynchronous RL (IMPALA, SEED~RL) decouples actor collection from learner
gradient updates to keep the GPU busy during rollouts.
Between callbacks, the system already achieves 860~fps and gradient updates
cost only 19~ms.  Async RL would give at most a 1.3$\times$ improvement
on the training loop and does nothing for callback overhead.

\textbf{Inference fix.} SB3's \texttt{model.predict} wrapper calls
\texttt{set\_training\_mode()} and allocates a new tensor on \textit{every}
invocation. For the callback's serial inference loop these overheads dominate.
The fix: extract the policy once, call \texttt{set\_training\_mode(False)}
once, pre-allocate a GPU buffer, and call \texttt{policy.\_predict} directly:

\begin{lstlisting}[language=Python]
# Once at callback setup (make_rl_alloc_cb):
policy = model.policy
policy.set_training_mode(False)
_obs_buf = torch.zeros(1, obs_dim, dtype=torch.float32, device=device)

# Per VM arrival (no alloc, no set_training_mode):
_obs_buf[0].copy_(torch.from_numpy(obs))   # zero-copy view + GPU transfer
with torch.no_grad():
    action = policy._predict(_obs_buf, deterministic=True).cpu().numpy()[0]
\end{lstlisting}

\subsection{Projected Callback Time After Both Fixes (Prediction)}

\begin{table}[h]
\centering
\begin{tabular}{lrrr}
\toprule
Phase & Before (run~v3) & After $D_j$ fix & After both fixes \\
\midrule
$D_j$/$S_j$ loop (O(\texttt{num\_mhd})/event) & $\approx$35 min & $<$1 s  & $<$1 s \\
Post-alloc tracking loop (O(\texttt{num\_mhd})/event) & included above & included above & $<$1 s \\
Inference (733k calls $\times$ 0.167~ms)               & included above & $\approx$122 s & $\approx$122 s \\
Numpy + purge overhead                                 & $\approx$2 min & $\approx$2 min & $\approx$10 s \\
\midrule
\textbf{Total per callback (predicted)} & \textbf{39 min} & \textbf{?} & \textbf{$\approx$2.2 min} \\
\bottomrule
\end{tabular}
\caption{Predicted callback cost. The $D_j$ fix was applied for run~v3 but the
post-allocation tracking loop was missed; run~v4 measures the actual result.}
\end{table}

% ─────────────────────────────────────────────────────────────────────────────
\section{Run~v4: Full 2M-Step Reward~A Run}
\label{sec:run4}
% ─────────────────────────────────────────────────────────────────────────────

\subsection{Configuration}

Run~v4 (\texttt{aug\_v3\_rewardA}) used the $D_j$ fix and the inference fix
but \textit{not} the post-allocation loop fix (discovered later):

\begin{lstlisting}[language=bash]
python -u scripts/train_rl.py \
  --run-id aug_v3_rewardA --reward-variant A \
  --augmentation --multi-trace \
  --aug-traces AMS20PrdApp19-tround ... \
  --n-envs 32 --total-timesteps 2000000 \
  --train-freq 8 --eval-freq 100000 \
  2>&1 | tee run4.txt
\end{lstlisting}

\subsection{Time Budget (Measured)}

\begin{table}[h]
\centering
\begin{tabular}{lrr}
\toprule
Phase & Duration & \% of total \\
\midrule
Startup (trace load + fork)        & 11 min 1 s  & 6.5\% \\
Actual training (env + grad)       & 13 min 35 s & 7.9\% \\
Callbacks (20 $\times$ 7.3~min)    & 146 min 28 s & \textbf{85.6\%} \\
\midrule
\textbf{Total}                     & \textbf{2 hr 50 min 39 s} & 100\% \\
\bottomrule
\end{tabular}
\caption{Exact time budget for run~v4 (2M timesteps).
20 callback windows averaged 480.2~s each (vs.\ 40.7~s of actual training).
The env+gradient system ran at $\approx$2,454~fps between callbacks,
consistent with the \texttt{train\_freq}=8 smoke test.
All overhead is in the callback.}
\end{table}

\subsection{Second O(\texttt{num\_mhd}) Bug: Post-Allocation Tracking Loop}

Despite the $D_j$/S$_j$ fix, each callback still took 7.3~min instead of
the predicted 2.5~min. Root cause: a second identical O(\texttt{num\_mhd})
loop in \texttt{pooling\_simulation} (evaluate.py, after \texttt{alloc\_fn}
returns):

\begin{lstlisting}[language=Python]
# BEFORE — O(11,904) per VM arrival event:
for j in range(num_mhd):
    if alloc_vec[j] > 0 and dealloc_time < pod_dur:
        mpd_vm_allocs[j].append((dealloc_time, float(alloc_vec[j])))
\end{lstlisting}

\texttt{alloc\_vec} is a 11,904-element array but only 8 entries
(the \texttt{mhd\_list} for the current host) are non-zero.
Iterating over all 11,904 elements to find 8 non-zero ones mirrors the
$D_j$/$S_j$ bug exactly.
With 733,400 total events across 10 simulation iterations:

\[
733{,}400 \times 11{,}904 = 8.73\,\text{B iterations} \approx 5\text{ min extra per callback}
\]

\textbf{Fix.} Iterate over \texttt{mhd\_list} directly (8 iterations, not 11,904):

\begin{lstlisting}[language=Python]
# AFTER — O(8) per VM arrival event:
for mhd in mhd_list:
    if alloc_vec[mhd] > 0 and dealloc_time < pod_dur:
        mpd_vm_allocs[mhd].append((dealloc_time, float(alloc_vec[mhd])))
\end{lstlisting}

\subsection{Predicted Time Budget After All Fixes (Run~v5)}

\begin{table}[h]
\centering
\begin{tabular}{lrr}
\toprule
Phase & Run~v4 (measured) & Run~v5 (predicted) \\
\midrule
Startup                  & 11 min  & 11 min \\
Training                 & 14 min  & 14 min \\
Callbacks (20 $\times$)  & 147 min (7.3~min each) & 44 min (2.2~min each) \\
\midrule
\textbf{Total}           & \textbf{2 hr 51 min}   & \textbf{$\approx$69 min} \\
\bottomrule
\end{tabular}
\caption{Predicted run~v5 time budget after fixing the post-allocation tracking loop.
The callback cost breakdown for run~v5: inference 122~s + D$_j$ $<$1~s +
post-alloc $<$1~s + numpy/purge 10~s $=$ 133~s $\approx$ 2.2~min per callback.}
\end{table}

\subsection{Learning Signal (Run~v4)}

\begin{table}[h]
\centering
\begin{tabular}{rrrr}
\toprule
Timesteps & \texttt{pooling\_savings\_mean} & \texttt{std} & \texttt{ent\_coef} \\
\midrule
100k  & 0.0137 &  0.0185 & 0.891 \\
300k  & 0.0138 &  0.0098 & 0.705 \\
500k  & 0.0136 &  0.0095 & 0.605 \\
700k  & 0.0122 &  0.0132 & 0.487 \\
900k  & 0.0121 &  0.0122 & 0.393 \\
1100k & 0.0148 &  0.0177 & 0.308 \\
1300k & 0.0133 &  0.0100 & 0.263 \\
1500k & 0.0133 &  0.0165 & 0.224 \\
1700k & 0.0041 &  0.0203 & 0.192 \\
1900k & 0.0018 &  0.0267 & 0.167 \\
2000k & 0.0142 &  0.0181 & 0.155 \\
\bottomrule
\end{tabular}
\caption{\texttt{pooling\_savings\_mean} = fraction by which Reward~A policy
outperforms greedy on LON23 eval trace. Policy achieves $\approx$1.4\%
savings early but does not improve monotonically. High variance relative
to mean suggests the policy has not converged. \texttt{ent\_coef} decays
steadily (0.891 $\to$ 0.155), indicating the policy is collapsing to a
near-deterministic strategy without finding a significantly better optimum
than greedy.}
\end{table}

\subsection{Execution Diagnoser}

A sidecar monitoring script (\texttt{scripts/monitor\_run.py}) was written to
diagnose future runs.  It samples GPU utilisation, CPU, and RAM every 10~s,
detects callback boundaries by tailing the run log, and produces a phase
breakdown and JSON sidecar on exit:

\begin{lstlisting}[language=bash]
# Terminal 2 (alongside training):
python scripts/monitor_run.py --log run5.txt --interval 10
\end{lstlisting}

% ─────────────────────────────────────────────────────────────────────────────
\section{Next Steps}
% ─────────────────────────────────────────────────────────────────────────────

\begin{enumerate}
  \item \textbf{Run~v5: Reward~A with post-allocation loop fix.}
        All three O(\texttt{num\_mhd}) bugs now fixed.
        Expected: $\approx$69~min total for 2M steps.
        Monitor with \texttt{scripts/monitor\_run.py} to verify callback drops
        to $\approx$2.2~min.

  \item \textbf{Convergence criterion.} No fixed step budget.
        Stop when \texttt{pooling\_savings\_mean} plateaus over
        3--5 consecutive eval windows.
        Target: $>0\%$ beats greedy; $>1.5\%$ is a strong result.
        Savings are currently flat ($\sim$1.2\% mean, high variance) ---
        2M steps may be insufficient; extend to 5M if needed.

  \item \textbf{Reward~B ablation} once Reward~A converges.
\end{enumerate}

\end{document}
